\section{Limitations and Future Directions}
\label{app:limitations}

The present work studies the shared-basis training-free transplant family (OMP~\citep{goddard2025tftransplant}; CLP~\citep{ostendorff2023efficient}, WECHSEL~\citep{minixhofer2021wechsel}, and FOCUS~\citep{dobler2023focus} in Appendix~\ref{sec:altops-results}) in the regime where shared-basis transplant is currently deployed (e.g., \texttt{mergekit-tokensurgeon}~\citep{mergekit2024}, deployed open-weight artifacts such as Arcee SuperNova-Medius). Appendix~\ref{app:nonlinear-shared-anchor} formalizes the conditions under which a nonlinear coefficient solver does not alter the conclusion. Principled algorithmic redesign of transplant operators, including recipient-side checks on synthesized rows, merge operators that explicitly handle appended rows absent from one merge input, and post-transplant behavioral validation, is the defense direction most directly motivated by our findings, since the LoRA suppression demonstrated in Appendix~\ref{app:lora_persistence} is bounded to prompts whose distribution matches the training corpus and does not cover the held-out out-of-distribution prompts we test, and is therefore not a sufficient mitigation against this attack family in our case study.
